%!TEX root = ../../csuthesis_research.tex

\chapter{各应用章节框架图与设计要点}

本章在第三章统一记号与递归闭式解的基础上，分别针对医学图像疾病分类、时间序列分类、语义分割以及多模态大语言模型专家路由四类任务给出对应的框架设计要点。各节均按照“场景特点与设计动机—关键模块与数据流—核心要点”的统一结构展开，旨在突出闭式解方法在不同任务形态下的统一性与差异性。

\section{医疗图像疾病分类持续学习框架}

医学图像疾病分类的持续学习面临隐私敏感、类别频繁扩展以及样本分布严重不平衡三方面挑战。一方面，病理切片与器官组织影像通常携带个人健康信息，难以直接保存历史样本进行回放；另一方面，新发疾病亚型与新采集模态的不断引入要求模型具备良好的长期可扩展性。为此，本节采用“预训练编码器加解析分类头递推”的解析持续学习策略，框架如图~\ref{fig:framework_ch3} 所示~\citep{li2026aidc}。

\begin{figure}[htbp]
    \centering
    \includegraphics[width=0.93\textwidth]{images/aidc-overview.pdf}
    \caption{闭式解持续学习方法在医学图像疾病分类中的三个主要组成部分：(a) 基于反向传播（Back-Propagation, BP）的训练，使用主干网络和 Softmax 分类器做多轮训练；(b) 基于解析学习的第一节段重对齐，把原分类器替换为解析分类器并与主干网络对齐，其中缓冲层是高维扩展层，通过最小二乘计算并保存特征的最优解；(c) 和 (d) 基于闭式解的持续学习，借助递归岭回归闭式解来更新解析分类头。}
    \label{fig:framework_ch3}
\end{figure}

\subsection{场景特点与设计动机}

医学影像数据的隐私属性使得回放型方法在临床部署中受到明显约束，而递归岭回归闭式解仅维护二阶统计量、不存储历史原始样本，恰好契合无回放需求。与此同时，临床场景中类别扩展呈高频低样本特征：每新增一种疾病亚型仅需一次解析更新即可纳入判别空间，无需对早期类别再次反向传播。此外，闭式解在理论上与联合训练等价，对类别不平衡所引发的判别边界偏移具备天然的鲁棒性，避免少数类被多数类淹没。

\subsection{关键模块与数据流}

\begin{itemize}
    \item \textbf{基础编码器}：在初始任务上以监督学习方式训练特征提取主干，获取可迁移的医学影像表征，并在第二阶段起冻结其参数。
    \item \textbf{解析对齐模块}：用解析分类头替换原 Softmax 分类头，借助最小二乘求解编码器输出到标签空间的最优投影，得到初始统计量 $\Psi_0$ 与 $\widehat{\Phi}_0$。
    \item \textbf{递归闭式解头}：当新阶段数据 $\mathcal{D}_t=\{\mathbf{X}_t,\mathbf{Y}_t\}$ 到达时，按第三章定理~\ref{thm:closed_form_recursive_report} 仅利用 $\Psi_{t-1}$、$\widehat{\Phi}_{t-1}$ 与当前样本递推得到 $\Psi_t$、$\widehat{\Phi}_t$。
    \item \textbf{评测与隐私接口}：以总体 Acc、平均 Forg 以及无样本回放约束为核心指标，输出对应阶段的判别边界与统计量快照，供下一阶段直接复用。
\end{itemize}

\subsection{该框架的核心要点}

第一，递归更新使历史知识以统计量形式被显式存储，而不是依赖梯度约束或样本回放进行隐式记忆，因而遗忘率更低；同时由于不再需要多轮反向传播，单阶段更新仅需一次特征遍历，更新效率显著提升。第二，在皮肤病、血液细胞与器官切片等多种医学影像类增量基准上验证了解析持续学习对隐私敏感与类别频繁扩展场景的适用性。第三，由于框架不引入任何额外可学习参数，整体结构清晰，易于在临床部署环境中与现有筛查流程集成。

\section{时间序列持续学习框架}

时序任务在持续学习中除灾难性遗忘之外，还须面对类内变化突出的难题：同一类别在不同个体、设备或环境下可能呈现差异较大的动态模式。为兼顾稳定性与可塑性，本节框架在闭式解之外引入多尺度特征融合与随机高维映射，并进一步推广出集成扩展形式，框架如图~\ref{fig:framework_ch4} 所示~\citep{li2025tsacl}。

\begin{figure}[htbp]
    \centering
    \includegraphics[width=0.98\textwidth]{images/tsacl-overview.pdf}
    \caption{所提方法总体框架。(a) 先在任务 1 上通过基于梯度下降的训练得到预训练编码器；(b) 接着在分类头前面引入带随机隐藏层模块，用以增强特征维度；(c) 最后在跨任务过程中使用递归岭回归学习机制。}
    \label{fig:framework_ch4}
\end{figure}

\subsection{场景特点与设计动机}

时间序列样本 $\mathbf{X}_t \in \mathbb{R}^{N_t \times c \times l}$ 由多通道传感器在时间维度上连续采样得到，其类内分布通常由多个子簇组成，单纯依赖小规模回放缓冲难以覆盖全部子模式。递归岭回归闭式解通过维护类内统计量而非具体样本来描述分布，因而天然适合刻画丰富的类内变化；新类别的引入也只需一次解析更新即可在判别空间中追加子空间，避免遗漏类内多样性。

\subsection{关键模块与数据流}

\begin{itemize}
    \item \textbf{一维卷积编码器}：在首任务上通过梯度下降训练时序主干，获取局部波形与全局语义的基础表征，训练完成后冻结。
    \item \textbf{多尺度特征拼接}：将编码器多层输出 $\{\mathbf{X}_t^{\mathrm{feat},k}\}_{k=1}^{K}$ 沿样本维度堆叠后再沿特征维度拼接，得到 $\mathbf{X}_t^{\mathrm{stack}}$，兼顾浅层局部波形与深层语义。
    \item \textbf{随机高维映射 (RHL)}：以随机初始化矩阵 $\mathbf{R}^E$ 将拼接特征非线性投影至高维空间 $\mathbf{X}_t^{E}=\mathrm{ReLU}(\mathbf{X}_t^{\mathrm{stack}}\mathbf{R}^E)$，提升解析分类头的线性可分性。
    \item \textbf{递归闭式解头}：直接套用第三章统一记号下的递推公式，仅借助 $\Psi_{t-1}$、$\widehat{\Phi}_{t-1}$ 与 $\mathbf{X}_t^{E}$ 完成增量更新。
    \item \textbf{集成扩展分支}：为同一编码器派生 $M$ 套随机映射矩阵 $\{\mathbf{R}^{E,i}\}_{i=1}^{M}$，并行维护 $M$ 个解析分类头，推理时按 Softmax 概率平均后再取最大类别。
\end{itemize}

\subsection{该框架的核心要点}

其一，多尺度特征融合使得在编码器冻结的前提下仍可显著提升对类内变化的可塑性，浅层特征保留波形细节、深层特征强化语义一致性。其二，集成扩展通过对随机隐藏层施加多种子初始化，构造出多个相互独立的高维判别空间，使预测在不同随机投影方向上得到互补：
\begin{equation}
\hat{y}=\arg\max_{c}\frac{1}{M}\sum_{i=1}^{M}P_i(c\mid \mathbf{x}^{\mathrm{test}}),
\end{equation}
计算开销仅以 $M$ 线性放大，并不引入额外反向传播。其三，由于递归更新结果在理论上与联合训练等价，编码器特征足够稳健时，解析分类头即可在无历史样本的条件下逼近“看过全部历史样本”的决策边界。

\section{语义分割持续学习框架}

语义分割将持续学习从图像级分类推广至像素级或点级的稠密预测，所面临的稳定性与可塑性矛盾、计算成本以及背景语义漂移问题均较图像分类更为严峻。本节框架在闭式解基础上引入伪标签机制，以统一方式处理二维图像与三维点云，框架如图~\ref{fig:framework_ch5} 所示~\citep{li2025cfsseg}。

\begin{figure}[htbp]
    \centering
    \includegraphics[width=0.98\textwidth]{images/cfsseg-overview.pdf}
    \caption{所提方法 CFSSeg 的总体框架。第 $t$ 步中，模型借助第 $t-1$ 步的模型通过伪标签机制生成伪标签，并与真实标签融合形成混合标签。模型继承第 $t-1$ 步学到的分类头 $\widehat{\Phi}_{t-1}$，再结合混合标签、提取出的高维特征矩阵以及第 $t-1$ 步的 $\Psi_{t-1}$。随后使用递归岭回归闭式解算法做更新，得到 $\widehat{\Phi}_{t}$ 与 $\Psi_{t}$。}
    \label{fig:framework_ch5}
\end{figure}

\subsection{场景特点与设计动机}

类增量分割中不相交与重叠设定下，背景标签可能携带旧类别甚至未来类别信息，直接训练易使旧类语义被背景类吸收，从而触发语义漂移。递归岭回归闭式解将每个位置级特征视为线性回归输入后，分类头更新即化为一次解析求解，避免了在数百万像素或点上反复执行反向传播；其递推形式在数学上等价于联合训练，使决策边界仅由累积统计量决定，对增量过程中的扰动具备良好鲁棒性。在此基础上引入伪标签可恢复被错误映射至背景的旧类监督信息，进而缓解语义漂移。

\subsection{关键模块与数据流}

\begin{itemize}
    \item \textbf{冻结分割编码器}：二维任务采用 DeepLabv3 加 ResNet-101~\citep{chen2017deeplabv3,he2016resnet}，三维任务采用 DGCNN~\citep{wang2019dgcnn}，初始任务后冻结主干。
    \item \textbf{随机高维映射层}：$\mathbf{X}_t^{E}=\mathrm{ReLU}(\mathbf{X}_t^{\mathrm{encoder}}\mathbf{R}^E)$，将位置级特征投影至高维空间，提升冻结编码器条件下的可塑性。
    \item \textbf{伪标签生成器}：以第 $t-1$ 阶段模型对当前数据中标注为背景的像素或点重新打标，构造伪监督矩阵。
    \item \textbf{递归闭式解头}：将编码后的位置级特征与伪标签融合后的标签矩阵代入第三章递推公式，得到新阶段的 $\Psi_t$ 与 $\widehat{\Phi}_t$。
    \item \textbf{推理输出}：在每个空间位置取分类头最大响应类别，形成完整的稠密预测结果。
\end{itemize}

\subsection{该框架的核心要点}

二维与三维伪标签策略在动机一致的前提下针对各自模态做了差异化设计。二维图像采用基于 Sigmoid 置信度的不确定性度量 $U^{i}=1-\sigma(\max_c q_{\theta_{t-1}}(i,c))$，当背景像素的旧模型置信度较高时以旧模型预测替换背景标签，否则保留原始背景标签。三维点云则进一步引入邻域加权熵度量，将近邻几何结构纳入考虑，以应对点云的局部稀疏性与几何连续性，并通过近邻投票得到点级伪标签 $\hat{Y}_{t-1}^{i,k'}$。

\section{多模态大语言模型专家路由框架}

多模态大语言模型在终身学习场景下通常采用“冻结主干加任务专家”的结构。新任务持续到达时，路由器自身亦面临持续学习挑战：一旦专家选择错误，即便专家本体完好，最终响应也会失败，故路由稳定性对整套系统起决定性作用。本节将专家选择视为解析分类问题，提出仅对路由头做闭式解递推、对任务专家以低秩适配训练的统一框架，结构如图~\ref{fig:framework_ch6} 所示~\citep{li2026cfrouter}。

\begin{figure}[htbp]
    \centering
    \includegraphics[width=0.95\textwidth]{images/cfrouter-overview.pdf}
    \caption{CF-Router 闭式解多模态专家路由框架概览。上图：终身学习场景，涵盖遥感、医疗、自动驾驶、科学以及金融等不同领域任务，它们按时间顺序依次到达，每个领域都有自己的视觉问答数据。下图：CF-Router 架构。多模态输入 $\mathbf{x}$（图像与问题）由冻结的多模态大语言模型主干 $f_{\theta}$ 处理，生成隐藏状态序列 $\mathbf{X}_{\mathrm{seq}} \in \mathbb{R}^{L \times D}$。平均池化将其压缩为固定大小特征 $\mathbf{X}_{\mathrm{feat}} \in \mathbb{R}^{D}$。训练期间增量累积自相关矩阵 $\mathbf{G}$ 与互相关矩阵 $\mathbf{Q}$，随后解析计算路由头 $\widehat{\Phi}_t = (\mathbf{G} + \gamma\mathbf{I})^{-1}\mathbf{Q}$，无需任何梯度下降。推理期间计算得分 $\mathbf{s} = \mathbf{X}_{\mathrm{feat}}^{\mathrm{test}}\widehat{\Phi}_t$，将得分最高的专家 $E_{k^*}$ 与冻结主干合并生成响应。}
    \label{fig:framework_ch6}
\end{figure}

\subsection{场景特点与设计动机}

多模态终身学习中，特征空间由图像、文本与跨模态对齐信息共同决定，简单的相似度度量难以稳定划分任务边界。冻结的多模态大语言模型主干 $f_{\theta}$ 在最后一层隐藏状态中已凝结深层多模态语义，可直接作为高判别力的路由特征；而专家编号作为多类标签后，路由问题即化为解析可解的多类判别问题，正好对应第三章统一框架。同时，闭式解递推使路由器无需重训亦无需回放历史多模态样本即可吸收新任务，与多模态数据规模庞大、回放代价高的现实约束相符。

\subsection{关键模块与数据流}

\begin{itemize}
    \item \textbf{冻结多模态主干} $f_{\theta}$：接收图文混合输入并输出最后一层隐藏状态序列 $\mathbf{X}_{\mathrm{seq}} \in \mathbb{R}^{L\times D}$，整个持续学习过程中保持不变。
    \item \textbf{序列池化}：以平均池化 $\mathbf{X}_{\mathrm{feat}}=\frac{1}{L}\sum_{i=1}^{L}\mathbf{X}_{\mathrm{seq},i}$ 得到固定维度语义特征，作为解析路由头的输入。
    \item \textbf{统计量缓冲区}：维护自相关矩阵 $\mathbf{G}^{t}=\sum_{i\le t}\mathbf{X}_i^{\top}\mathbf{X}_i$ 与互相关矩阵 $\mathbf{Q}^{t}=\sum_{i\le t}\mathbf{X}_i^{\top}\mathbf{Y}_i$，按批次累积更新。
    \item \textbf{解析路由头}：在阶段末按 $\widehat{\Phi}_t=(\mathbf{G}^{t}+\gamma\mathbf{I})^{-1}\mathbf{Q}^{t}$ 一次性解析求解，与第三章 $\widehat{\Phi}_t=\Psi_t\mathbf{Q}_t$ 严格等价。
    \item \textbf{LoRA 任务专家}：对每个任务 $t$ 初始化轻量级低秩适配模块 $E_t$ 并通过梯度下降训练，训练完成后冻结。
    \item \textbf{推理路由}：先按 $\mathbf{s}=\mathbf{X}_{\mathrm{feat}}^{\mathrm{test}}\widehat{\Phi}_t$ 计算得分，再取 $k^{*}=\arg\max_k s_k$ 激活对应专家与冻结主干合并生成响应。
\end{itemize}

\subsection{该框架的核心要点}

第一，专家选择头的闭式解更新使路由器在新任务到达时仅需累积一次统计量并执行一次解析求解，避免了反向传播带来的梯度遗忘问题，且 $\mathbf{G}^{t}$ 与 $\Psi_t$ 之间满足 $\mathbf{G}^{t}=\Psi_t^{-1}-\gamma\mathbf{I}$，与第三章递推完全等价。第二，整套框架既是解析的、也是无回放的：除统计量缓冲区外，不存储任何历史多模态样本，与多模态终身学习对存储和隐私的双重约束相符。第三，由于路由稳定性直接决定专家激活是否正确，本节同时关注准确率与后向迁移指标 BWT，借此衡量长期任务序列下的路由稳定性与遗忘抑制能力。

\section{四类任务的统一性回顾}

尽管医学影像疾病分类、时序分类、语义分割与多模态专家路由四类任务在输入模态、输出粒度与监督形式上差异显著，但其更新方式均可归结为第三章给出的递归岭回归闭式解：在任务 $t$ 到达时仅需累积当前阶段的 $\mathbf{X}_t^{\top}\mathbf{X}_t$ 与 $\mathbf{X}_t^{\top}\mathbf{Y}_t$，再借助 $\Psi_{t-1}$、$\widehat{\Phi}_{t-1}$ 完成一次解析更新即可。各任务的差异主要体现在编码器选择、特征构造以及监督信号的具体形式上，而递推主体保持一致。



\section{本章小结}

本章对四类应用场景的框架设计要点进行了系统整理，明确了它们在第三章统一记号下的对应关系：编码器负责模态相关的特征抽取，高维映射负责提升冻结主干条件下的线性可分性，递归闭式解头负责无回放增量更新，任务适配模块则负责伪标签生成、集成扩展或路由解析等差异化处理。下一章将围绕这一统一框架展开实验设计，依次给出各任务的基准数据集、对比基线、关键超参数以及上下界设定，以验证闭式解持续学习方法在不同应用形态下的普适性与有效性。
